% ============================================================
% 修士論文アブストラクト本文（2頁）
% PDF上では自動的に2段組へ流れます．\switchcolumnは使いません．
% \AbstractPageBreakより前が1頁目，後が2頁目です．この命令は削除しない
% でください．内容の区切りに応じた移動は可能ですが，正確に2頁であることと
% 両頁の密度を確認してください．架空の数値や未実施の結果は書かないこと．
% 以下の7節は推奨構成例です．節名・節数は変更できますが，目的，差分，
% 方法，評価，限界，結論が追える構成にしてください．
% 文献は\cite{key}で引用します．\nociteと\printbibliographyは書きません．
% 参考文献一覧はabstract.texが独立アブストラクトへ出力します．
% 修士論文本体はこのファイルを読み込まず，本文用文献だけを別に管理します．
% ============================================================

\AbstractSection{1．緒言}
% 研究背景，未解決課題，学術的・技術的な必要性を記載する．
% 最後に，研究目的と検証する問いを明示する．

\AbstractSection{2．関連研究}
% 研究群を比較軸で整理し，本研究が埋める差分を明確にする．
% 引用する場合はreferences.bibへ登録し，本文中で\cite{key}を使う．

\AbstractSection{3．問題設定}
% 入力，出力，仮定，記号，目的関数，適用範囲を定義する．

\AbstractSection{4．提案手法}
% 全体構成，主要な新規要素，学習・推論手順を再現可能な粒度で記載する．
% 図と番号付き数式を置いた場合は，本文中で図\ref{fig:...}，式\eqref{eq:...}のように参照する．
% 図題は図の下へ置く．

% ここが物理的な1頁目と2頁目の境界です．
\AbstractPageBreak

\AbstractSection{5．実験・評価}
% データ，分割，ベースライン，評価指標，実装条件，主要結果を記載する．
% アブレーション，感度分析，統計的ばらつきも必要に応じて示す．
% 表を置いた場合は本文中で表\ref{tab:...}のように参照し，表題は表の上へ置く．

\AbstractSection{6．考察}
% 主張を結果へ対応づけ，限界，失敗条件，外的妥当性を述べる．
% 観測事実と推測を分ける．

\AbstractSection{7．結言}
% 目的に対する答え，主要な貢献，残された課題を簡潔にまとめる．
